% =============================================================
%  CS102-F26 Digital Design and Computer Architecture I
%  Lecture 04 -- Boolean Functions (Unit 4)
%  Follows Sections 2.3 and 2.3.1 of A Student's Computer Architecture Book I.
%
%  Build the slides and the print handout:  npm run materials -- lecture-04
% =============================================================
\documentclass{wosm-slides}
\usepackage{cs102-lectures}[2026/10/19]

\title[Lecture 04 \textperiodcentered\ Boolean Functions]{Boolean Functions}
\subtitle{Lecture 04: truth tables, and the canonical sum-of-products and
  product-of-sums forms}

% ---------- Helpers for this lecture --------------------------
% Proof steps: an equation and the law that justifies it.
\newenvironment{derivation}
  {\begin{tabular}{@{}r@{\;}l@{\qquad}l@{}}}
  {\end{tabular}}
\newcommand*\by[1]{\wosmlabel{#1}}

% A two-level circuit: one gate per term, all feeding one output gate.
% \twolevel{<first gate>}{<output gate>}{<terms>}{<output gate inputs>}
\newcommand*\twolevel[4]{%
  \begin{tikzpicture}[gates]
    \foreach \t [count=\i from 0] in {#3} {
      \node[#1, inputs=nnn] (g\i) at (0,-\i*0.8) {};
      \foreach \k in {1,2,3} \draw[wire] (g\i.input \k) -- ++(-0.25,0);
      \node[left, font=\scriptsize] at ($(g\i.input 2)+(-0.25,0)$) {$\t$};
      \xdef\twolevellast{\i}
    }
    \node[#2, inputs=#4, logic gate input sep=0.2cm] (out)
      at (2.5,{-\twolevellast*0.4}) {};
    \foreach \i [evaluate=\i as \j using int(\i+1)] in {0,...,\twolevellast}
      \draw[wire] (g\i.output) -- ++(0.3,0) -- ($(out.input \j)+(-0.3,0)$)
        -- (out.input \j);
    \draw[wire] (out.output) -- ++(0.4,0) node[right] {$f$};
  \end{tikzpicture}}

% =============================================================
\begin{document}

\maketitle

\begin{frame}{Today}
  \tableofcontents
\end{frame}

% =============================================================
\section{Boolean functions}

\begin{frame}{What is a Boolean function?}
  \begin{definition}[Boolean function]
    A \emph{Boolean function} of $n$ variables is a function
    \[ f : \{0, 1\}^n \to \{0, 1\} \]
    built from the Boolean operations. Put simply: anything that takes one or more
    0/1 inputs and gives one 0/1 output.
  \end{definition}

  \medskip
  \begin{itemize}
    \item $f(x, y, z) = xy + z$ is a Boolean function of three variables.
    \item So is any lookup table that lists every 3-bit input on the left and one
      bit on the right: the table defines the function.
  \end{itemize}
\end{frame}

\begin{frame}{The truth table}
  \begin{columns}[T, totalwidth=\textwidth]
    \column{0.42\textwidth}
      \centering
      \begin{tabular}{@{}c|ccc|c|c@{}}
        row & $x$ & $y$ & $z$ & $xy$ & $f = xy + z$ \\
        \midrule
        0 & 0 & 0 & 0 & 0 & 0 \\
        1 & 0 & 0 & 1 & 0 & 1 \\
        2 & 0 & 1 & 0 & 0 & 0 \\
        3 & 0 & 1 & 1 & 0 & 1 \\
        4 & 1 & 0 & 0 & 0 & 0 \\
        5 & 1 & 0 & 1 & 0 & 1 \\
        6 & 1 & 1 & 0 & 1 & 1 \\
        7 & 1 & 1 & 1 & 1 & 1 \\
      \end{tabular}
    \column{0.52\textwidth}
      \begin{itemize}
        \item One row for every input: $2^3 = 8$ rows.
        \item Rows are listed in counting order. The \alert{row number} is the input
          $xyz$ read as a binary number: row 6 is $x = 1$, $y = 1$, $z = 0$.
        \item To fill it in, evaluate the expression row by row; helper columns,
          like $xy$, keep it tidy.
      \end{itemize}
  \end{columns}
\end{frame}

\begin{frame}{What a truth table is, and is not}
  A truth table is
  \begin{itemize}
    \item \textbf{tabular}, not algebraic;
    \item \textbf{declarative}, not imperative: it says \emph{which} output belongs to
      each input, not \emph{how} to compute it;
    \item \textbf{exhaustive}: every input must be there. With 10 inputs that is
      $2^{10} = 1024$ rows.
  \end{itemize}

  \bigskip
  To build a circuit we need an algebraic form that says how to compute $f$. The
  next two sections give two standard ones, the \alert{canonical} forms, that can
  be read straight off any truth table.
\end{frame}

\begin{frame}{How many Boolean functions are there?}
  A function of $n$ variables is a truth table with $2^n$ rows, and each row's
  output is 0 or 1. So there are $2^{2^n}$ different functions.

  \medskip
  \begin{columns}[T, totalwidth=\textwidth]
    \column{0.42\textwidth}
      \centering
      \begin{tabular}{@{}rrr@{}}
        \toprule
        $n$ & rows & functions \\
        \midrule
        1 & 2 & 4 \\
        2 & 4 & 16 \\
        3 & 8 & 256 \\
        4 & 16 & 65,536 \\
        \bottomrule
      \end{tabular}
    \column{0.52\textwidth}
      \centering
      \wosmlabel{The four functions of one variable}\par\smallskip
      \begin{tabular}{@{}c|cccc@{}}
        $x$ & $0$ & $x'$ & $x$ & $1$ \\
        \midrule
        0 & 0 & 1 & 0 & 1 \\
        1 & 0 & 0 & 1 & 1 \\
      \end{tabular}
  \end{columns}

  \bigskip
  Many of the 16 functions of two variables have names: AND, OR, XOR, NAND, NOR,
  \dots{} They become the logic gates of Unit~6.
\end{frame}

% =============================================================
\section{Sum of products}

\begin{frame}{Detecting one row}
  Build a function that is 1 for the input on \alert{one} row, and 0 everywhere else.
  Take row 2: $x = 0$, $y = 1$, $z = 0$.

  \[
    r_2(x, y, z) =
      \begin{cases}
        1 & \text{if } x = 0 \text{ and } y = 1 \text{ and } z = 0, \\
        0 & \text{otherwise.}
      \end{cases}
  \]

  \begin{itemize}
    \item A product is 1 only when \emph{every} factor is 1.
    \item $x'$ is 1 exactly when $x = 0$; $y$ when $y = 1$; $z'$ when $z = 0$.
    \item So $r_2(x, y, z) = x'yz'$.
  \end{itemize}

  \medskip
  Each variable appears once: as it is where its bit is 1, complemented where it
  is 0.
\end{frame}

\begin{frame}{Minterms}
  \begin{columns}[T, totalwidth=\textwidth]
    \column{0.55\textwidth}
      \begin{definition}[Minterm]
        A \emph{minterm} of $n$ variables is a product in which every variable
        appears once, plain or complemented. The minterm $m_i$ is 1 on row $i$
        only.
      \end{definition}

      \smallskip
      A row detector like $r_2$ is a minterm: $r_2 = m_2$.
    \column{0.42\textwidth}
      \centering
      \small
      \begin{tabular}{@{}cccc|l@{}}
        row & $x$ & $y$ & $z$ & minterm \\
        \midrule
        0 & 0 & 0 & 0 & $m_0 = x'y'z'$ \\
        1 & 0 & 0 & 1 & $m_1 = x'y'z$ \\
        2 & 0 & 1 & 0 & $m_2 = x'yz'$ \\
        3 & 0 & 1 & 1 & $m_3 = x'yz$ \\
        4 & 1 & 0 & 0 & $m_4 = xy'z'$ \\
        5 & 1 & 0 & 1 & $m_5 = xy'z$ \\
        6 & 1 & 1 & 0 & $m_6 = xyz'$ \\
        7 & 1 & 1 & 1 & $m_7 = xyz$ \\
      \end{tabular}
  \end{columns}
\end{frame}

\begin{frame}{The canonical sum of products}
  $f = xy + z$ is 1 on rows 1, 3, 5, 6 and 7. So $f$ is 1 exactly when \emph{one of}
  their detectors is 1: OR them together.

  \[
    f = m_1 + m_3 + m_5 + m_6 + m_7
      = x'y'z + x'yz + xy'z + xyz' + xyz
  \]

  \begin{block}{Sum of minterms}
    Every Boolean function is the sum of the minterms of the rows where it is 1.
    In short: $f(x, y, z) = \sum m(1, 3, 5, 6, 7)$.
  \end{block}

  For a given order of variables this form is unique: two expressions are the same
  function exactly when they have the same minterms.
\end{frame}

\begin{frame}{As a circuit: AND--OR}
  \begin{columns}[c, totalwidth=\textwidth]
    \column{0.5\textwidth}
      \centering
      \twolevel{and gate}{or gate}{x'y'z, x'yz, xy'z, xyz', xyz}{nnnnn}
    \column{0.46\textwidth}
      \begin{itemize}
        \item One AND gate per minterm: its output is 1 only when it sees its row.
        \item One OR gate collects them: $f$ is 1 when any row detector is 1.
        \item Complemented inputs come from NOT gates (not drawn).
      \end{itemize}
      \medskip
      \wosmlabel{The book's Figure 2.1 shows the same circuit in Logisim-evolution.
      Gates in detail: Unit 6.}
  \end{columns}
\end{frame}

% =============================================================
\section{Product of sums}

\begin{frame}{Detecting the zeros}
  Instead, detect the rows where $f = 0$: rows 0, 2 and 4.

  \medskip
  \begin{itemize}
    \item $m_0 + m_2 + m_4$ is 1 exactly where $f$ is 0: it is $f'$.
    \item So $f$ is its complement. Apply De Morgan (Lecture 03), twice:
  \end{itemize}

  \medskip
  \small
  \begin{derivation}
    $f$ & $= (m_0 + m_2 + m_4)' = m_0' \cdot m_2' \cdot m_4'$ & \by{De Morgan} \\
        & $= (x'y'z')' \cdot (x'yz')' \cdot (xy'z')'$ & \\
        & $= (x + y + z)(x + y' + z)(x' + y + z)$ & \by{De Morgan, involution} \\
  \end{derivation}
\end{frame}

\begin{frame}{Maxterms}
  \begin{columns}[T, totalwidth=\textwidth]
    \column{0.55\textwidth}
      \begin{definition}[Maxterm]
        A \emph{maxterm} of $n$ variables is a sum in which every variable appears
        once, plain or complemented. The maxterm $M_i = m_i'$ is 0 on row $i$
        only.
      \end{definition}

      \smallskip
      Each variable appears as it is where its bit is \alert{0}, complemented where
      it is 1: the opposite of a minterm.
    \column{0.42\textwidth}
      \centering
      \small
      \begin{tabular}{@{}cccc|l@{}}
        row & $x$ & $y$ & $z$ & maxterm \\
        \midrule
        0 & 0 & 0 & 0 & $M_0 = x + y + z$ \\
        1 & 0 & 0 & 1 & $M_1 = x + y + z'$ \\
        2 & 0 & 1 & 0 & $M_2 = x + y' + z$ \\
        3 & 0 & 1 & 1 & $M_3 = x + y' + z'$ \\
        4 & 1 & 0 & 0 & $M_4 = x' + y + z$ \\
        5 & 1 & 0 & 1 & $M_5 = x' + y + z'$ \\
        6 & 1 & 1 & 0 & $M_6 = x' + y' + z$ \\
        7 & 1 & 1 & 1 & $M_7 = x' + y' + z'$ \\
      \end{tabular}
  \end{columns}
\end{frame}

\begin{frame}{The canonical product of sums}
  \begin{columns}[c, totalwidth=\textwidth]
    \column{0.5\textwidth}
      \begin{block}{Product of maxterms}
        Every Boolean function is the product of the maxterms of the rows where it
        is 0:
        \[ f(x, y, z) = \prod M(0, 2, 4). \]
      \end{block}
      \begin{itemize}
        \item A product is 0 when any factor is 0, and $M_i$ is 0 only on row $i$.
        \item So $f$ is 0 on rows 0, 2 and 4, and 1 everywhere else.
      \end{itemize}
    \column{0.46\textwidth}
      \centering
      \twolevel{or gate}{and gate}{x+y+z, x+y'+z, x'+y+z}{nnn}\par
      \medskip
      \wosmlabel{OR--AND circuit; compare Figure 2.2 of the book.}
  \end{columns}
\end{frame}

\begin{frame}{Two forms, one function}
  \[
    f = \textstyle\sum m(1, 3, 5, 6, 7) = \prod M(0, 2, 4)
  \]

  \begin{itemize}
    \item The two index lists together contain every row exactly once: the minterms
      are the 1s, the maxterms are the 0s.
    \item The complement swaps the roles:
      $f' = \sum m(0, 2, 4) = \prod M(1, 3, 5, 6, 7)$.
    \item The sum of products has one term per 1, the product of sums one term per
      0. Use the form with fewer terms.
  \end{itemize}

  \begin{warning}[Common mistake]
    In a minterm a 0 bit gives a complemented variable; in a maxterm a 0 bit gives
    a plain one. Row 4 ($x = 1$, $y = 0$, $z = 0$): $m_4 = xy'z'$, but
    $M_4 = x' + y + z$.
  \end{warning}
\end{frame}

% =============================================================
\section{Working with canonical forms}

\begin{frame}{From an expression to the sum of products}
  Multiply each product by $(v + v') = 1$ for every variable $v$ it is missing,
  expand, and drop repeated terms.

  \medskip
  \small
  \begin{derivation}
    $xy + z$ & $= xy(z + z') + (x + x')(y + y')\,z$ & \by{complement, identity} \\
             & $= xyz + xyz' + x'y'z + x'yz + xy'z + xyz$ & \by{distributivity} \\
             & $= xyz + xyz' + x'y'z + x'yz + xy'z$ & \by{idempotence} \\
             & $= \sum m(1, 3, 5, 6, 7)$ & \\
  \end{derivation}

  \medskip
  \normalsize
  The same five minterms as the truth table gave us.
\end{frame}

\begin{frame}{From an expression to the product of sums}
  First turn the expression into a product of sums; then add $v v' = 0$ for every
  missing variable $v$, and split with the second distributive law.

  \medskip
  \small
  \begin{derivation}
    $xy + z$ & $= (x + z)(y + z)$ & \by{distributivity} \\
    $x + z$ & $= x + yy' + z = (x + y + z)(x + y' + z)$ & \by{complement, distributivity} \\
    $y + z$ & $= xx' + y + z = (x + y + z)(x' + y + z)$ & \by{complement, distributivity} \\
  \end{derivation}

  \medskip
  \normalsize
  Multiply, and keep $(x + y + z)$ once (idempotence):
  \[ xy + z = (x + y + z)(x + y' + z)(x' + y + z) = \textstyle\prod M(0, 2, 4). \]
\end{frame}

\begin{frame}{Try it: the majority function}
  $\mathit{maj}(x, y, z)$ is 1 when at least two of its three inputs are 1.

  \medskip
  \begin{columns}[T, totalwidth=\textwidth]
    \column{0.3\textwidth}
      \centering
      \small
      \begin{tabular}{@{}ccc|c@{}}
        $x$ & $y$ & $z$ & $\mathit{maj}$ \\
        \midrule
        0 & 0 & 0 & \uncover<2->{0} \\
        0 & 0 & 1 & \uncover<2->{0} \\
        0 & 1 & 0 & \uncover<2->{0} \\
        0 & 1 & 1 & \uncover<2->{1} \\
        1 & 0 & 0 & \uncover<2->{0} \\
        1 & 0 & 1 & \uncover<2->{1} \\
        1 & 1 & 0 & \uncover<2->{1} \\
        1 & 1 & 1 & \uncover<2->{1} \\
      \end{tabular}
    \column{0.64\textwidth}
      \begin{enumerate}
        \item Fill in the truth table.
        \item Write the canonical sum of products.\par
          \uncover<3->{\small $\sum m(3, 5, 6, 7) = x'yz + xy'z + xyz' + xyz$}
        \item Write the canonical product of sums.\par
          \uncover<4->{\small $\prod M(0, 1, 2, 4)
            = (x + y + z)(x + y + z')(x + y' + z)(x' + y + z)$}
      \end{enumerate}
      \medskip
      \uncover<4->{\wosmlabel{The majority of three bits is the carry when you add
      them (Unit 16).}}
  \end{columns}
\end{frame}

% =============================================================
\section{Summary}

\begin{frame}{Summary}
  \begin{itemize}
    \item A Boolean function maps $n$ input bits to one output bit. Its truth table
      has $2^n$ rows; there are $2^{2^n}$ such functions.
    \item A minterm $m_i$ is 1 on row $i$ only; a maxterm $M_i = m_i'$ is 0 on row
      $i$ only.
    \item Canonical sum of products: the minterms of the 1-rows, $\sum m(\dots)$,
      built as AND--OR.
    \item Canonical product of sums: the maxterms of the 0-rows, $\prod M(\dots)$,
      built as OR--AND.
    \item Both forms are unique, but long. Next lecture: making them short.
  \end{itemize}
\end{frame}

\begin{frame}{Before the next lecture}
  \begin{itemize}
    \item Read Section~2.3 of the book, up to the end of \emph{Canonical
      Representations}.
    \item Write $\sum m(0, 3, 5, 6)$ as a product of maxterms, and write out both
      forms in full.
    \item Expand $x + y'z$ into its canonical sum of products. How many minterms
      does it have?
  \end{itemize}

  \begin{block}{Next lecture \unitref{5}}
    Simplifying Boolean functions with Karnaugh maps.
  \end{block}
\end{frame}

\end{document}
